\section{Biometría hemática}

La biometría hemática (BH), o hemograma, es un estudio sanguineo realizado para la identificación de enfermedades o fallas metabólicas. De manera general, requiere de entre 4 y 8 horas de ayuno, y su lugar de punción es en las venas. Ahora bien, existen diferencias entre la realización de una biometría hemática básica y una total.

     \begin{table}
\centering
\caption{Comparativa entre biometría hemática básica y total}
\label{tab:comparar_biometrías}
\begin{threeparttable}
{\setstretch{1.75}
\begin{tabular}{ccc}\hline

    \textbf{Valor} & \textbf{Básica} & \textbf{Total} \\ \hline
    
    GR\tnote{1} & Conteo total de GR &  Conteo de total de GR, de hemoglobina y de hematocrito \\    
    GB\tnote{2} & Conteo total de GB & Conteo diferencial (contabiliza cada tipo de GB) \\
    Plaquetas & Cantidad y tamaño & Cantidad y tamaño\\
    
    \hline
\end{tabular}
}
\begin{tablenotes}
    \item[1] Glóbulos rojos
    \item[2] Glóbulos blancos
\end{tablenotes}
\end{threeparttable}
\end{table}

Como se puede observar en la \autoref{tab:comparar_biometrías}

\section{Bioquímica sanguinea}

Mientras que una biometría hemática se enfoca en analizar las células de la sangre, una bioquímica sanguinea analiza, como su nombre lo dice, los compuestos bioquímicos presentes en ella, esto con objetivos similares a la biometría: descubrir la existencia de alguna falla inmune o metabólica en el paciente, así como el correcto funcionamiento de los órganos del cuerpo

Este procedimiento requiere de un ayuno de 8-12 horas por parte del paciente.

Los compuestos más importantes a analizar durante una bioquímica sanguinea, también conocidos como perfiles, son:

\begin{itemize}
    \item Glucosa
    \item Urea
    \item Ácido úrico
    \item Creatinina
    \item Colesterol
    \item Triglicéridos
\end{itemize}

Además de estos perfiles, se pueden analizar hasta otros 45 elementos mediante las pruebas de sangre.